% !PRINTED 335
\weberpage{335}
%<<CONT
in which $f_k,\Delta_k,J_k$ denote the derivatives with respect to $x_k$. For the canonical form we then have
\weq{11}{K(y) = r^9 K.}
To find $K(y)$, we need only form the functional determinant from the left-hand sides of the formulae (7), (8), (9), from which, by application of simple determinant theorems, we obtain (Bd. I, \S.\,22, \S.\,27):
\[
K(y) = -h^3(h^3+8m^3)^3 \left|\begin{array}{ccc} 1, & y_1^3, & y_1^6 \\ 1, & y_2^3, & y_2^6 \\ 1, & y_3^3, & y_3^6 \end{array}\right|,
\]
and consequently
\weq{12}{r^9 K = h^3(h^3+8m^3)^3\,(y_1^3-y_2^3)(y_1^3-y_3^3)(y_2^3-y_3^3).}
It is still to be remarked concerning these formulae that the factor $h^3+8m^3$ occurring in them cannot vanish if the curve $f$ has no singular point. For a singular point is determined by the three equations:
\[
h\,y_1^2 + 2m\,y_2y_3 = 0,\quad h\,y_2^2 + 2m\,y_3y_1 = 0,\quad h\,y_3^2 + 2m\,y_1y_2 = 0,
\]
and these three equations are compatible with one another when $h^3+8m^3=0$. (If, for example, $h=-2m$, they are satisfied when $y_1=y_2=y_3$.)

The coefficient $h$ can have any prescribed value different from zero for a given form $f$, and only when this value is given is the problem determined. For the sake of simplicity we will now set $h=1$, and obtain from (7) to (9) the result:

If one sets
\weq{13}{\begin{aligned}
P &= \frac{(2+m^3)m^3f^2+(2-5m^3)mr^2f\Delta+3m^2r^4\Delta^2-r^6J}{(1+8m^3)^2}, \\
Q &= \frac{(1+2m^3)f-6mr^2\Delta}{1+8m^3}, \\
R &= \frac{m^2f+r^2\Delta}{1+8m^3},
\end{aligned}}
then $y_1^3,y_2^3,y_3^3$ are the roots of the cubic equation
\weq{14}{u^3 - Qu^2 + Pu - R^3 = 0,}
and the discriminant of this cubic equation is:
\weq{15}{D_1 = \frac{r^{18}K^2}{(1+8m^3)^6}.}
